\documentclass[12pt]{article}
\usepackage[utf8]{inputenc}
\usepackage[T1]{fontenc}
\usepackage{hyperref}
\title{Nugunu Language Resources for K-12 Curriculum Development}
\author{Compiled by Claude Sonnet 5.5 (claude-sonnet-5-5) for Andr\'as Kornai}
\date{}
\begin{document}
\maketitle

\section{Introduction}
Nugunu (Gunu, Ngunu, Yambasa) is a Bantu language (Zone A, Sanaga group, A.622) of the Centre region of Cameroon, in Mbam-et-Inoubou. It has Northern and Southern dialects. A figure of 35,000 speakers dates from 1987. The searches found little material, and only that is listed here.

\section{Primary Educational Resources}
No primers, textbooks or school programmes were found in the searches.

\section{Linguistic and Reference Materials}
\begin{itemize}
\item Wikipedia entry for the language: \url{https://www.wikipedia.com/wiki/ISO_639:yas}
\item Language Description Heritage listing, which includes a 1991 study of tone in the verb phrase (Patman) and a 1980 manuscript on participants in narrative discourse (Robinson): \url{https://ldh.clld.org/languages/yas}
\item PHOIBLE phonological inventory: \url{https://phoible.org/languages/yas}
\item Webonary dictionary and introduction site (Gunu): \url{https://www.webonary.work/gunu/overview/introduction/?lang=en}. I did not check its contents beyond the title.
\end{itemize}

\section{Religious and Community Resources}
\begin{itemize}
\item Scripture Earth listing of scripture resources: \url{https://www.scriptureearth.org/00eng.php?iso=yas}. A search summary mentions the Gospel of Luke on Bible.is and a YouVersion text.
\item Joshua Project, Yambasa, Nu Gunu: \url{https://joshuaproject.net/people_groups/15971}
\end{itemize}

\section{Digital and Audio Resources}
Global Recordings Network has a page for the language, but a search summary says it lists no recordings: \url{https://globalrecordings.net/en/language/yas}

\section{Curriculum Development Context}
Nugunu is spoken by the whole community and is stable, but no school use was found. A curriculum would need a standard orthography and reading materials. The Webonary dictionary site and the scripture texts could be a starting point for vocabulary and reading texts. Local partners would need to be identified, and none were found in the searches.

\end{document}